\documentclass[11pt]{article}

\usepackage[margin=1in]{geometry}
\usepackage[T1]{fontenc}
\usepackage[utf8]{inputenc}
\usepackage{lmodern}
\usepackage{microtype}
\usepackage{amsmath, amssymb, amsfonts, mathtools}
\usepackage{bm}
\usepackage{booktabs}
\usepackage{array}
\usepackage{enumitem}
\usepackage{xcolor}
\usepackage{hyperref}
\usepackage{listings}

\hypersetup{
  colorlinks=true,
  linkcolor=blue!50!black,
  urlcolor=blue!50!black,
  citecolor=blue!50!black,
  pdftitle={Exam 01 Solution - Machine Learning Theory},
  pdfauthor={Codex}
}

\lstdefinestyle{pythonstyle}{
  language=Python,
  basicstyle=\ttfamily\small,
  keywordstyle=\color{blue!60!black},
  commentstyle=\color{green!40!black},
  stringstyle=\color{red!50!black},
  showstringspaces=false,
  breaklines=true,
  frame=single,
  columns=fullflexible
}

\setlist[itemize]{topsep=4pt, itemsep=2pt, leftmargin=1.5em}

\newcommand{\R}{\mathbb{R}}
\newcommand{\Hspace}{\mathcal{H}}
\newcommand{\E}{\mathbb{E}}
\newcommand{\Var}{\operatorname{Var}}
\newcommand{\MMD}{\operatorname{MMD}}
\newcommand{\norm}[1]{\left\lVert #1 \right\rVert}
\newcommand{\inner}[2]{\left\langle #1, #2 \right\rangle}
\newcommand{\transpose}{^{\top}}

\title{\textbf{Exam 01 Solution}\\[4pt]\large Machine Learning Theory}
\author{Universidad Nacional de Colombia\\Course deliverable in English}
\date{Term 2026-1}

\begin{document}

\maketitle
\thispagestyle{empty}

\begin{center}
\begin{tabular}{>{\bfseries}rl}
Source statement: & \texttt{exams/01/statement/main.pdf} \\
Code verification: & \texttt{exams/01/solutions/code/solve\_exam\_01.py} \\
Numeric outputs: & \texttt{exams/01/solutions/artifacts/numeric\_results.json}
\end{tabular}
\end{center}

\vspace{1em}

\noindent
This document provides a polished English solution to the first exam. The original statement numbers the SVD/kernel-trick sub-question as \texttt{1.3} again; here it is relabeled as \textbf{1.4} for clarity.

\tableofcontents
\newpage

\section{Regression in Feature Space and Kernel Formulation}

\subsection{Regularized Least Squares in Feature Space}

We observe i.i.d. pairs
\[
\{(x_n, t_n)\}_{n=1}^N,
\qquad
x_n \in \R^P,
\qquad
t_n \in \R,
\]
and we assume the model
\[
t_n = \inner{w}{\phi(x_n)}_{\Hspace} + \eta_n,
\qquad
\eta_n \sim \mathcal{N}(0, \sigma_\eta^2).
\]

The unknown parameter is the weight vector
\[
w \in \Hspace.
\]

The regularized least-squares problem in feature space is
\begin{equation}
\widehat{w}
=
\arg\min_{w \in \Hspace}
\left[
\frac{1}{2}
\sum_{n=1}^N
\Bigl(t_n - \inner{w}{\phi(x_n)}_{\Hspace}\Bigr)^2
+
\frac{\lambda}{2}\norm{w}_{\Hspace}^2
\right],
\qquad
\lambda \ge 0.
\label{eq:rrls-feature}
\end{equation}

In compact notation, if $t = [t_1,\dots,t_N]\transpose$ and $(\Phi w)_n = \inner{w}{\phi(x_n)}_{\Hspace}$, then
\[
\widehat{w}
=
\arg\min_{w \in \Hspace}
\left[
\frac{1}{2}\norm{t - \Phi w}_2^2
+
\frac{\lambda}{2}\norm{w}_{\Hspace}^2
\right].
\]

The regularization parameter $\lambda$ plays four roles:
\begin{itemize}
\item it penalizes large feature-space norms $\norm{w}_{\Hspace}$;
\item it improves conditioning when $\Phi\transpose \Phi$ is singular or nearly singular;
\item it controls the bias-variance tradeoff;
\item it discourages fitting the observation noise too aggressively.
\end{itemize}

\subsection{Closed-Form Solution in Matrix Notation}

Assume first that the feature space admits a finite coordinate representation of size $M$. Then
\[
\Phi
=
\begin{bmatrix}
\phi(x_1)\transpose \\
\vdots \\
\phi(x_N)\transpose
\end{bmatrix}
\in \R^{N \times M},
\qquad
t =
\begin{bmatrix}
t_1 \\
\vdots \\
t_N
\end{bmatrix}
\in \R^N,
\qquad
w \in \R^M.
\]

The objective in \eqref{eq:rrls-feature} becomes
\[
J(w)
=
\frac{1}{2}(t-\Phi w)\transpose (t-\Phi w)
+
\frac{\lambda}{2} w\transpose w.
\]

Expanding the quadratic term gives
\[
J(w)
=
\frac{1}{2}
\left(
t\transpose t
- 2 t\transpose \Phi w
+ w\transpose \Phi\transpose \Phi w
\right)
+
\frac{\lambda}{2} w\transpose w.
\]

Differentiating with respect to $w$,
\[
\nabla_w J(w)
=
-\Phi\transpose t + \Phi\transpose \Phi w + \lambda w.
\]

The first-order optimality condition is
\[
-\Phi\transpose t + \Phi\transpose \Phi w + \lambda w = 0,
\]
which is equivalent to
\begin{equation}
(\Phi\transpose \Phi + \lambda I_M)w = \Phi\transpose t.
\label{eq:normal-equation}
\end{equation}

Hence the estimator is
\begin{equation}
\widehat{w}
=
(\Phi\transpose \Phi + \lambda I_M)^{-1}\Phi\transpose t.
\label{eq:primal-solution}
\end{equation}

In operator notation, valid even when $\Hspace$ is infinite-dimensional,
\[
\widehat{w}
=
(\Phi^\ast \Phi + \lambda I_{\Hspace})^{-1}\Phi^\ast t.
\]

\subsection{Polynomial Feature Map, Design Matrix, and Python Implementation}

Consider the polynomial feature map
\[
\phi(x)
=
\begin{bmatrix}
1 \\
x \\
x^2
\end{bmatrix}.
\]

For a dataset $\{(x_n, t_n)\}_{n=1}^N$, the associated design matrix is
\[
\Phi
=
\begin{bmatrix}
1 & x_1 & x_1^2 \\
1 & x_2 & x_2^2 \\
\vdots & \vdots & \vdots \\
1 & x_N & x_N^2
\end{bmatrix}
\in \R^{N \times 3}.
\]

The implementation in the repository uses the exact formula from \eqref{eq:primal-solution}. The core code is:

\begin{lstlisting}[style=pythonstyle]
def build_polynomial_design_matrix(x_samples: np.ndarray) -> np.ndarray:
    x_vector = np.asarray(x_samples, dtype=np.float64).reshape(-1)
    return np.column_stack((np.ones_like(x_vector), x_vector, x_vector**2))

def solve_regularized_least_squares(
    design_matrix: np.ndarray,
    targets: np.ndarray,
    lambda_reg: float,
) -> np.ndarray:
    gram = design_matrix.T @ design_matrix
    rhs = design_matrix.T @ targets
    system = gram + lambda_reg * np.eye(design_matrix.shape[1], dtype=np.float64)
    return np.linalg.solve(system, rhs)
\end{lstlisting}

\paragraph{Numerical example.}
Take
\[
x = [-2,-1,0,1,2]\transpose,
\qquad
t = 1 + 0.5x + x^2
= [4,1.5,1,2.5,6]\transpose.
\]

Then
\[
\Phi =
\begin{bmatrix}
1 & -2 & 4 \\
1 & -1 & 1 \\
1 & 0 & 0 \\
1 & 1 & 1 \\
1 & 2 & 4
\end{bmatrix}.
\]

Using $\lambda = 0.1$, the script returns
\[
\widehat{w}
\approx
\begin{bmatrix}
0.9674 \\
0.4950 \\
1.0066
\end{bmatrix}.
\]

The fitted outputs are
\[
\widehat{t}
\approx
[4.0038,\ 1.4790,\ 0.9674,\ 2.4691,\ 5.9840]\transpose.
\]

These values are very close to the generating polynomial coefficients $(1,0.5,1)$. The small discrepancy is entirely expected: the penalty term shrinks the solution slightly toward the origin. The script also verifies that the primal and dual solutions coincide up to machine precision; the maximum absolute difference is approximately $5.55 \times 10^{-17}$.

\subsection{Kernel Trick from the SVD of the Design Matrix}

Let the thin singular value decomposition of the design matrix be
\[
\Phi = U_r \Sigma_r V_r\transpose,
\]
where
\[
U_r \in \R^{N \times r},
\qquad
\Sigma_r \in \R^{r \times r},
\qquad
V_r \in \R^{M \times r},
\qquad
r = \operatorname{rank}(\Phi).
\]

Then
\[
\Phi\transpose \Phi = V_r \Sigma_r^2 V_r\transpose,
\qquad
\Phi \Phi\transpose = U_r \Sigma_r^2 U_r\transpose.
\]

The second matrix is the Gram matrix
\[
K = \Phi \Phi\transpose,
\qquad
K_{ij} = \inner{\phi(x_i)}{\phi(x_j)}_{\Hspace} = k(x_i,x_j).
\]

From the primal estimator,
\[
\widehat{w}
=
(\Phi\transpose \Phi + \lambda I)^{-1}\Phi\transpose t,
\]
substituting the SVD gives
\[
\widehat{w}
=
\left(V_r \Sigma_r^2 V_r\transpose + \lambda I\right)^{-1}
V_r \Sigma_r U_r\transpose t
=
V_r (\Sigma_r^2 + \lambda I_r)^{-1}\Sigma_r U_r\transpose t.
\]

Now define the dual coefficient vector
\[
\alpha = (K + \lambda I_N)^{-1}t.
\]

Using the SVD form of $K$,
\[
\alpha
=
U_r (\Sigma_r^2 + \lambda I_r)^{-1} U_r\transpose t
+
\text{components in } \ker(K).
\]

Multiplying by $\Phi\transpose$,
\[
\Phi\transpose \alpha
=
V_r \Sigma_r U_r\transpose
U_r (\Sigma_r^2 + \lambda I_r)^{-1} U_r\transpose t
=
V_r (\Sigma_r^2 + \lambda I_r)^{-1}\Sigma_r U_r\transpose t
=
\widehat{w}.
\]

Therefore
\[
\widehat{w} = \Phi\transpose \alpha,
\qquad
\alpha = (K+\lambda I_N)^{-1}t.
\]

For a new input $x$, the prediction is
\[
\widehat{f}(x)
=
\inner{\widehat{w}}{\phi(x)}_{\Hspace}
=
\left\langle \sum_{n=1}^N \alpha_n \phi(x_n), \phi(x) \right\rangle_{\Hspace}
=
\sum_{n=1}^N \alpha_n k(x_n,x).
\]

Equivalently,
\[
\widehat{f}(x)
=
k_x\transpose (K + \lambda I_N)^{-1} t,
\qquad
k_x =
\begin{bmatrix}
k(x_1,x) \\
\vdots \\
k(x_N,x)
\end{bmatrix}.
\]

This is the kernel trick: explicit feature coordinates are unnecessary, because the full computation is expressed in terms of kernel evaluations.

\section{RKHS, Mean Embeddings, and Discrepancy Between Distributions}

\subsection{Empirical Mean Embedding and Kernel Evaluations}

Let $k : \mathcal{X} \times \mathcal{X} \to \R$ be a positive-definite kernel with associated RKHS $\mathcal{H}_k$. Given a sample $\{x_n\}_{n=1}^N$, the empirical mean embedding is
\[
\widehat{\mu}_X
=
\frac{1}{N}\sum_{n=1}^N \phi(x_n)
=
\frac{1}{N}\sum_{n=1}^N k(x_n,\cdot)
\in \mathcal{H}_k.
\]

For any query point $z \in \mathcal{X}$, the reproducing property yields
\[
\inner{\widehat{\mu}_X}{k(z,\cdot)}_{\mathcal{H}_k}
=
\frac{1}{N}\sum_{n=1}^N
\inner{k(x_n,\cdot)}{k(z,\cdot)}_{\mathcal{H}_k}
=
\frac{1}{N}\sum_{n=1}^N k(x_n,z).
\]

So the embedding can be evaluated using kernel values only.

\paragraph{Linear kernel.}
For
\[
k(x_i,x_j) = x_i\transpose x_j,
\qquad
x_i,x_j \in \R^P,
\]
the RKHS is simply $\R^P$ with the standard inner product, and $\phi(x)=x$. Thus
\[
\widehat{\mu}_X
=
\frac{1}{N}\sum_{n=1}^N x_n
=
\bar{x}
\in \R^P.
\]

Its evaluation at $z \in \R^P$ is
\[
\inner{\widehat{\mu}_X}{z}
=
\bar{x}\transpose z
=
\frac{1}{N}\sum_{n=1}^N x_n\transpose z
=
\frac{1}{N}\sum_{n=1}^N k(x_n,z).
\]

Dimensions:
\begin{itemize}
\item $x_n, z \in \R^P$,
\item $\widehat{\mu}_X \in \R^P$,
\item the evaluation $\bar{x}\transpose z$ is a scalar.
\end{itemize}

\paragraph{Gaussian RBF kernel.}
For
\[
k(x_i,x_j) = \exp\bigl(-\gamma \norm{x_i-x_j}_2^2\bigr),
\qquad
\gamma > 0,
\]
the RKHS is generally infinite-dimensional, but the mean embedding is still
\[
\widehat{\mu}_X = \frac{1}{N}\sum_{n=1}^N k(x_n,\cdot),
\]
and its evaluation at $z$ is
\[
\inner{\widehat{\mu}_X}{k(z,\cdot)}_{\mathcal{H}_k}
=
\frac{1}{N}\sum_{n=1}^N \exp\bigl(-\gamma \norm{x_n-z}_2^2\bigr).
\]

Dimensions:
\begin{itemize}
\item $x_n, z \in \R^P$,
\item $\widehat{\mu}_X \in \mathcal{H}_k$,
\item the evaluation is a scalar,
\item in practice we work with the vector $k_z = [k(x_1,z), \dots, k(x_N,z)]\transpose \in \R^N$.
\end{itemize}

\paragraph{Interpretation.}
The linear-kernel embedding is exactly the ordinary sample mean in the original space. The RBF embedding instead summarizes the sample in a nonlinear feature space and therefore captures local similarity structure rather than only first-order Euclidean coordinates.

\paragraph{Numerical check.}
Using the sample
\[
\{(0,1), (1,0), (2,1)\}
\qquad
\text{and}
\qquad
z=(1,2),
\]
the script returns:
\begin{itemize}
\item linear sample mean: $(1,\tfrac{2}{3})$;
\item linear evaluation via $\bar{x}\transpose z$: $2.3333$;
\item direct linear kernel average: $2.3333$;
\item RBF evaluation with $\gamma = 0.5$: $0.2904$.
\end{itemize}

\subsection{Raw, Centered, and Standardized Moments}

The distinction between raw, centered, and standardized moments is conceptually identical in the original space and in the feature space; only the object whose moments are taken changes.

\paragraph{Original space.}
Let $X \in \R^P$ be a random vector with mean $\mu = \E[X]$ and marginal standard deviations collected in a diagonal matrix $S$.

The $m$th raw moment is
\[
\E[X^{\otimes m}],
\]
so it measures the moment relative to the origin.

The $m$th centered moment is
\[
\E[(X-\mu)^{\otimes m}],
\]
so it measures variation around the mean.

The $m$th standardized moment is
\[
\E\left[\left(S^{-1}(X-\mu)\right)^{\otimes m}\right],
\]
which removes both location and scale.

\paragraph{Kernel-induced space.}
Replace $X$ by the feature-space random element $\phi(X) \in \mathcal{H}_k$, and denote its mean by
\[
\mu_\phi = \E[\phi(X)].
\]

Then:
\begin{itemize}
\item raw moments are taken relative to the origin of $\mathcal{H}_k$;
\item centered moments are taken relative to $\mu_\phi$;
\item standardized moments remove both the mean effect and the scale effect.
\end{itemize}

Along any fixed direction $u \in \mathcal{H}_k$,
\[
\text{raw: }
\E\bigl[\inner{\phi(X)}{u}^m\bigr],
\]
\[
\text{centered: }
\E\bigl[\inner{\phi(X)-\mu_\phi}{u}^m\bigr],
\]
\[
\text{standardized: }
\E\left[
\left(
\frac{\inner{\phi(X)-\mu_\phi}{u}}
{\sqrt{\Var(\inner{\phi(X)}{u})}}
\right)^m
\right].
\]

Therefore:
\begin{itemize}
\item raw moments depend on the origin;
\item centered moments remove the mean;
\item standardized moments remove mean and scale simultaneously;
\item in RKHS this is especially important because nonlinear embeddings can alter norms and geometry substantially.
\end{itemize}

\subsection{RKHS Discrepancy Between Two Distributions and the MNIST Experiment}

Let $X \sim P$ and $Y \sim Q$, with embeddings
\[
\mu_P = \E[k(X,\cdot)],
\qquad
\mu_Q = \E[k(Y,\cdot)].
\]

The squared RKHS discrepancy is
\[
\norm{\mu_P - \mu_Q}_{\mathcal{H}_k}^2.
\]

Expanding the norm gives
\[
\norm{\mu_P - \mu_Q}_{\mathcal{H}_k}^2
=
\inner{\mu_P}{\mu_P}_{\mathcal{H}_k}
+
\inner{\mu_Q}{\mu_Q}_{\mathcal{H}_k}
- 2\inner{\mu_P}{\mu_Q}_{\mathcal{H}_k}.
\]

Using the reproducing property,
\[
\inner{\mu_P}{\mu_P}_{\mathcal{H}_k} = \E[k(X,X')],
\qquad
\inner{\mu_Q}{\mu_Q}_{\mathcal{H}_k} = \E[k(Y,Y')],
\qquad
\inner{\mu_P}{\mu_Q}_{\mathcal{H}_k} = \E[k(X,Y)].
\]

Hence
\begin{equation}
\MMD^2(P,Q)
=
\E[k(X,X')]
+
\E[k(Y,Y')]
- 2\E[k(X,Y)].
\label{eq:mmd-population}
\end{equation}

\paragraph{Empirical matrix form.}
Let
\[
X_m = \{x_1,\dots,x_m\},
\qquad
Y_n = \{y_1,\dots,y_n\}.
\]

Define
\[
(K_{XX})_{ij} = k(x_i,x_j),
\qquad
(K_{YY})_{ij} = k(y_i,y_j),
\qquad
(K_{XY})_{ij} = k(x_i,y_j).
\]

Then the biased empirical estimator is
\begin{equation}
\widehat{\MMD}_{\mathrm{b}}^2
=
\frac{1}{m^2}\bm{1}_m\transpose K_{XX}\bm{1}_m
+
\frac{1}{n^2}\bm{1}_n\transpose K_{YY}\bm{1}_n
- \frac{2}{mn}\bm{1}_m\transpose K_{XY}\bm{1}_n.
\label{eq:mmd-empirical}
\end{equation}

This is the matrix form requested by the exam.

\paragraph{MNIST implementation.}
The repository script downloads MNIST, selects a balanced subset of 100 training images per digit, uses an RBF kernel, and computes pairwise class discrepancies using \eqref{eq:mmd-empirical}. The bandwidth is selected by a median heuristic on a pooled subset, which in this run gives
\[
\gamma \approx 0.0097933.
\]

The smallest discrepancies are shown in Table~\ref{tab:smallest-mmd}, and the largest in Table~\ref{tab:largest-mmd}.

\begin{table}[h]
\centering
\caption{Five smallest empirical RKHS discrepancies on MNIST.}
\label{tab:smallest-mmd}
\begin{tabular}{ccr}
\toprule
Digit A & Digit B & Biased $\widehat{\MMD}^2$ \\
\midrule
4 & 9 & 0.0689 \\
3 & 5 & 0.0876 \\
5 & 8 & 0.0884 \\
7 & 9 & 0.0996 \\
8 & 9 & 0.1234 \\
\bottomrule
\end{tabular}
\end{table}

\begin{table}[h]
\centering
\caption{Five largest empirical RKHS discrepancies on MNIST.}
\label{tab:largest-mmd}
\begin{tabular}{ccr}
\toprule
Digit A & Digit B & Biased $\widehat{\MMD}^2$ \\
\midrule
0 & 1 & 0.4687 \\
1 & 7 & 0.3420 \\
1 & 4 & 0.3416 \\
1 & 6 & 0.3293 \\
1 & 9 & 0.3264 \\
\bottomrule
\end{tabular}
\end{table}

\paragraph{Interpretation.}
A small RKHS discrepancy means that the two empirical class distributions induce similar kernel mean embeddings. A large discrepancy means that the classes are well separated in the geometry induced by the kernel. In this run:
\begin{itemize}
\item the pair $4$ vs.\ $9$ is the closest, which is plausible because many handwritten instances share similar topological and stroke patterns;
\item the pair $0$ vs.\ $1$ is the most separated, which is also intuitive because one class is loop-dominated while the other is largely stroke-dominated;
\item several of the largest discrepancies involve the digit $1$, suggesting that under the chosen kernel it is especially distinctive relative to other classes.
\end{itemize}

These values are empirical and therefore depend on the sample size, the sampling seed, and the kernel bandwidth.

\section{Matrix Form of Kernels and Induced Geometry}

\subsection{Squared Euclidean Distances in Matrix Form}

Take two sample rows $x_i, x_j \in \R^P$. Then
\[
\norm{x_i-x_j}_2^2
=
(x_i-x_j)\transpose (x_i-x_j).
\]

Expanding the product gives
\[
\norm{x_i-x_j}_2^2
=
x_i\transpose x_i
+
x_j\transpose x_j
- 2x_i\transpose x_j.
\]

Now define the data matrix
\[
X =
\begin{bmatrix}
x_1\transpose \\
\vdots \\
x_N\transpose
\end{bmatrix}
\in \R^{N \times P},
\qquad
G = XX\transpose \in \R^{N \times N}.
\]

Let the vector of squared norms be
\[
s =
\begin{bmatrix}
\norm{x_1}_2^2 \\
\vdots \\
\norm{x_N}_2^2
\end{bmatrix}
=
\operatorname{diag}(G)
\in \R^N,
\]
and let $\bm{1} \in \R^N$ denote the all-ones vector. Then the full matrix of squared distances is
\begin{equation}
D
=
s \bm{1}\transpose + \bm{1} s\transpose - 2XX\transpose.
\label{eq:distance-matrix}
\end{equation}

Indeed, the $(i,j)$ entry of the right-hand side is
\[
s_i + s_j - 2x_i\transpose x_j
=
\norm{x_i}_2^2 + \norm{x_j}_2^2 - 2x_i\transpose x_j
=
\norm{x_i-x_j}_2^2.
\]

\subsection{Gaussian Kernel Matrix: Diagonal, Symmetry, and Positive Semidefiniteness}

The Gaussian (RBF) kernel matrix is
\[
K_{ij} = \exp\bigl(-\gamma \norm{x_i-x_j}_2^2\bigr),
\qquad
\gamma > 0.
\]

\paragraph{Diagonal entries.}
If $i=j$, then $\norm{x_i-x_i}_2^2 = 0$, so
\[
K_{ii} = \exp(0) = 1.
\]
Hence all diagonal entries are identical:
\[
K_{11} = K_{22} = \cdots = K_{NN} = 1.
\]

Geometrically,
\[
K_{ii}
=
\inner{\phi(x_i)}{\phi(x_i)}_{\Hspace}
=
\norm{\phi(x_i)}_{\Hspace}^2,
\]
so every transformed point has the same RKHS norm:
\[
\norm{\phi(x_i)}_{\Hspace} = 1.
\]
Therefore all embedded samples lie on the unit sphere of the induced feature space.

\paragraph{Symmetry.}
Since Euclidean distance is symmetric,
\[
\norm{x_i-x_j}_2^2 = \norm{x_j-x_i}_2^2,
\]
it follows immediately that
\[
K_{ij} = K_{ji}.
\]
Therefore
\[
K = K\transpose.
\]

\paragraph{Positive semidefiniteness.}
The Gaussian kernel is positive definite. Therefore, for every coefficient vector $c \in \R^N$,
\[
\sum_{i=1}^N \sum_{j=1}^N c_i c_j k(x_i,x_j) \ge 0,
\]
which is equivalent to
\[
c\transpose K c \ge 0.
\]
Hence $K$ is positive semidefinite.

Because the kernel is positive definite, there exists a feature map $\phi$ such that
\[
K_{ij} = k(x_i,x_j) = \inner{\phi(x_i)}{\phi(x_j)}_{\Hspace}.
\]
So the Gaussian kernel matrix is precisely a Gram matrix in the induced feature space.

\paragraph{Numerical sanity check.}
The verification script confirms the theory:
\begin{itemize}
\item the maximum error between the direct distance computation and the matrix identity \eqref{eq:distance-matrix} is approximately $5.68 \times 10^{-14}$;
\item the maximum deviation of the RBF diagonal from $1$ is approximately $5.55 \times 10^{-16}$;
\item the symmetry error is zero up to machine precision;
\item on the tested subset, the minimum eigenvalue is approximately $0.0777$, which is consistent with positive semidefiniteness.
\end{itemize}

\section{Conclusion}

This exam links four mutually reinforcing viewpoints:
\begin{itemize}
\item primal regularized least squares in feature space;
\item dual kernel regression through the Gram matrix;
\item RKHS mean embeddings and MMD as distribution-level summaries;
\item Gaussian kernels as distance-based Gram matrices with a clear induced geometry.
\end{itemize}

The unifying principle is that once feature-space inner products are replaced by kernel evaluations, the entire analysis can be carried out without explicitly constructing the feature map.

\end{document}
